\section{Praktikum}
\label{sec:praktikum-bab5}

\begin{praktikumbox}
\textbf{Judul:} Pengamatan register dan \eng{flag} dengan
\eng{Single Step} pada Emu8086.\\[0.3em]
\textbf{Sub-CPMK:} 1.2 \quad\textbullet\quad
\textbf{Alokasi waktu:} 100 menit \quad\textbullet\quad
\textbf{Bentuk:} individu
\end{praktikumbox}

\subsection{Tujuan Praktikum}

Setelah menyelesaikan praktikum ini, mahasiswa mampu mengamati isi
setiap register 8086, membaca keadaan \eng{flag} setelah operasi
aritmatika, dan menjelaskan akibat perubahan nilai register secara
manual terhadap jalannya program.

\subsection{Alat dan Bahan}

\begin{itemize}
    \item Komputer dengan Emu8086 terpasang.
    \item Berkas \file{prediksi-flag.asm} dan \file{pesan-flag.asm}
          dari direktori \file{code/asm/}.
    \item Lembar kerja pengamatan register dan \eng{flag}.
\end{itemize}

\subsection{Langkah Kerja}

\begin{enumerate}
    \item \textbf{Menyiapkan program pengamatan.}
    \begin{enumerate}
        \item Buka Emu8086, lalu buka berkas
              \file{prediksi-flag.asm}.
        \item Tekan \file{Compile and Emulate}, kemudian tahan pada
              jendela emulator tanpa menjalankan program.
        \item Pastikan jendela \file{Registers} dan jendela
              \file{Flags} terlihat. Bila belum, aktifkan melalui
              menu \file{View}.
    \end{enumerate}

    \item \textbf{Mengamati register dengan \eng{Single Step}.}
    \begin{enumerate}
        \item Tekan tombol \file{Single Step} satu kali untuk
              menjalankan instruksi pertama.
        \item Catat nilai \reg{AX}, \reg{BX}, \reg{CX}, \reg{DX},
              \reg{DS}, \reg{SP}, dan \reg{IP} pada
              Tabel~\ref{tab:pengamatan-register} setiap kali satu
              instruksi selesai.
        \item Berhenti pada instruksi \lstasm|ADD AX, BX| dan catat
              keadaan \eng{flag} sebelum instruksi dijalankan.
        \item Jalankan instruksi \lstasm|ADD AX, BX| satu langkah,
              lalu catat keadaan \eng{flag} sesudahnya.
        \item Lanjutkan sampai program selesai dan catat keadaan
              \eng{flag} pada akhir kasus pertama dan kasus kedua.
    \end{enumerate}

    \item \textbf{Mengisi tabel pengamatan \eng{flag}.}
    \begin{enumerate}
        \item Isi kolom prediksi pada Tabel~\ref{tab:prediksi-flag}
              \emph{sebelum} menjalankan kasus yang bersangkutan.
        \item Bandingkan prediksi dengan hasil pengamatan.
        \item Untuk setiap prediksi yang salah, tuliskan penjelasan
              mengapa \eng{flag} tersebut bernilai demikian.
    \end{enumerate}

    \item \textbf{Mengubah nilai register secara manual.}
    \begin{enumerate}
        \item Pada jendela \file{Registers}, klik ganda nilai
              \reg{AX}, ubah menjadi \hex{FFFF}, lalu tekan
              \file{Enter}.
        \item Jalankan satu instruksi berikutnya dan amati pengaruh
              nilai baru tersebut terhadap hasil dan \eng{flag}.
        \item Ubah nilai \reg{AL} saja menjadi \hex{00} dan amati
              bahwa \reg{AH} tidak ikut berubah.
        \item Ulangi langkah serupa untuk \reg{BX} dan catat
              perbedaannya.
    \end{enumerate}

    \item \textbf{Menguji pengaruh perubahan register kendali.}
    \begin{enumerate}
        \item Jalankan berkas \file{pesan-flag.asm} dengan
              \eng{Single Step}.
        \item Pada jendela \file{Flags}, ubah nilai \flag{CF} secara
              manual dari 1 menjadi 0, lalu lanjutkan eksekusi.
        \item Amati pesan mana yang tercetak dan jelaskan sebabnya
              berdasarkan urutan pemeriksaan \eng{flag}.
        \item Kembalikan nilai \flag{CF} dan ulangi percobaan untuk
              \flag{ZF}.
    \end{enumerate}

    \item \textbf{Mengamati register \reg{SP} dan \reg{SS}.}
    \begin{enumerate}
        \item Tuliskan potongan program singkat berisi
              \lstasm|MOV AX, 1234h| lalu \lstasm|PUSH AX|.
        \item Jalankan program dengan \eng{Single Step} dan catat
              nilai \reg{SP} sebelum dan sesudah \asm{PUSH}.
        \item Jelaskan mengapa \reg{SP} berubah sebesar 2, bukan 1.
    \end{enumerate}
\end{enumerate}

\subsection{Tabel Pengamatan}

\begin{table}[H]
    \centering
    \caption{Lembar pengamatan isi register pada praktikum bab ini.}
    \label{tab:pengamatan-register}
    \small
    \begin{tabular}{@{}lp{2.0cm}p{2.0cm}p{2.0cm}p{2.0cm}@{}}
        \toprule
        \textbf{Instruksi} & \reg{AX} & \reg{BX} & \reg{SP} &
        \reg{IP} \\
        \midrule
        \lstasm|MOV AX, @DATA| & \dots & \dots & \dots & \dots \\
        \addlinespace
        \lstasm|MOV DS, AX| & \dots & \dots & \dots & \dots \\
        \addlinespace
        \lstasm|MOV AX, 007Fh| & \dots & \dots & \dots & \dots \\
        \addlinespace
        \lstasm|MOV BX, 0001h| & \dots & \dots & \dots & \dots \\
        \addlinespace
        \lstasm|ADD AX, BX| & \dots & \dots & \dots & \dots \\
        \addlinespace
        \lstasm|MOV hasil1, AX| & \dots & \dots & \dots & \dots \\
        \bottomrule
    \end{tabular}
\end{table}

\subsection{Hasil yang Dikumpulkan}

\begin{enumerate}
    \item Tabel~\ref{tab:pengamatan-register} yang telah terisi
          lengkap.
    \item Tabel~\ref{tab:prediksi-flag} yang telah terisi prediksi
          dan hasil pengamatan.
    \item Tangkapan layar jendela \file{Registers} dan
          \file{Flags} pada saat \lstasm|ADD AX, BX| selesai
          dijalankan.
    \item Catatan hasil percobaan mengubah \flag{CF} dan \flag{ZF}
          secara manual beserta penjelasan sebabnya.
\end{enumerate}

\begin{asesmenbox}
\textbf{Asesmen praktikum 3 (bobot 5\% dari komponen praktikum).}

\begin{tabular}{@{}p{6.2cm}cp{4.6cm}@{}}
    \toprule
    \textbf{Indikator} & \textbf{Bobot} & \textbf{Kriteria} \\
    \midrule
    Ketepatan pengamatan register & 30\% &
    Seluruh nilai register tercatat benar pada setiap langkah \\
    \addlinespace
    Ketepatan pengamatan \eng{flag} & 30\% &
    Nilai \flag{CF}, \flag{OF}, \flag{SF}, \flag{ZF} sesuai hasil
    pengamatan \\
    \addlinespace
    Analisis perubahan manual & 25\% &
    Penjelasan sebab perubahan hasil dan \eng{flag} logis dan tepat \\
    \addlinespace
    Kerapian dokumentasi & 15\% &
    Lembar kerja lengkap, rapi, dan dapat diulang \\
    \bottomrule
\end{tabular}
\end{asesmenbox}
